\BourbakiUnnumberedChapter{Bibliography}{Bibliography}

\begin{enumerate}
\def\labelenumi{\arabic{enumi}.}
\item
  O. NEUGEBAUER, Vorlesungen über Geschichte der antiken Mathematik, Bd. I: Vorgriechische Mathematik, Berlin (Springer), 1934.
\item
  Euclidis Elementa, 5 vol., ed.~J. L. Heiberg, Lipsiae (Teubner), 1883-88.
\end{enumerate}

2 bis. T. L. HEATH, The thirteen books of Euclid's Elements\ldots, 3 vol., Cambridge, 1908.

\begin{enumerate}
\def\labelenumi{\arabic{enumi}.}
\setcounter{enumi}{2}
\item
  H. CARDANO, Opera, Lyon, 1663.
\item
  R. BOMBELLI, L'Algebra, Bologne (G. Rossi), 1572.
\item
  FRANCISCI VIETAE, Opera mathematica\ldots, Lugduni Batavorum (Elzevir), 1646.
\item
  A. GIRARD, Invention nouvelle en Algèbre, Amsterdam, 1629.
\item
  R. DESCARTES, Geometria, Latin transl. by Fr.~van Schooten, 2 ed., 2 vol., Amsterdam (Elzevir), 1659-61.
\item
  G. W. LEIBNIZ, Mathematische Schriften, ed.~C. I. Gerhardt, 7 vol., Berlin-Halle (Ascher-Schmidt), 1849-63.
\item
  Der Briefwechsel von Gottfried Wilhelm Leibniz mit Mathematikern, herausg. von C. I. Gerhardt, vol.~I, Berlin (Mayer und Müller), 1899.
\item
  L. EULER, Opera Omnia (1), vol.~VI, Berlin-Leipzig (Teubner), 1921: a) De Formis Radicum Aequationum\ldots, p.~1-19; b) De Resolutione Aequationum cujusvis gradus, p.~170-196.
\item
  J.-L. LAGRANGE, Œuvres, t. III, Paris (Gauthier-Villars), 1869 : a) Réflexions sur la résolution algébrique des équations, p.~205-421 ; b) Sur la forme des racines imaginaires des équations, p.~479.
\item
  A. VANDERMONDE, Mémoire sur la résolution des équations, Hist. de l'Acad. royale des sciences, année 1771, Paris (1774), p.~365-416.
\item
  C. F. GAUSS, Werke, vol.~I-X, Göttingen, 1863-1923.
\item
  N. H. ABEL, Œuvres, 2 vol., ed.~Sylow et Lie, Christiania, 1881.
\item
  E. GALOIS, Œuvres mathématiques, Paris (Gauthier-Villars), 1897.
\item
  A.-L. CAUCHY, Œuvres complètes (1), t. X, Paris (Gauthier-Villars), 1897, p.~312 et 351.
\item
  J. LIOUVILLE, Sur des classes très étendues de quantités dont la valeur n'est ni algébrique, ni même réductible à des irrationnelles algébriques, Journ. de Math. (1), t. XVI (1851), p.~133.
\item
  R. DEDEKIND, Gesammelte mathematische Werke, 3 vol., Braunschweig (Vieweg), 1932.
\item
  L. KRONECKER: a) Grundzüge einer arithmetischen Theorie der algebraischen Grössen, Crelle's J., XCII (1882), p.~1-122 (= Werke, vol.~II, Leipzig (Teubner), 1897, p.~245-387); b) Ein Fundamentalsatz der allgemeinen Arithmetik, Crelle's J., C (1887), p.~490-510 (= Werke, vol.~III \(_{1}\) , Leipzig (Teubner), 1899, p.~211-240).
\item
  H. WEBER, Untersuchungen über die allgemeinen Grundlagen der Galoisschen Gleichungstheorie, Math. Ann., vol.~XLIII (1893), p.~521-544.
\item
  G. VERONESE, Fondamenti di geometria, Padova, 1891.
\item
  K. HENSEL, Theorie der algebraischen Zahlen, Leipzig-Berlin (Teubner), 1908.
\item
  E. STEINITZ, Algebraische Theorie der Körper, Crelle's J., vol.~CXXXVII (1910), p.~167-309.
\item
  W. KRULL, Galoissche Theorie der unendlichen algebraischen Erweiterungen, Math. Ann., vol.~C (1928), p.~687.
\item
  E. ARTIN, Galois Theory\ldots, Ann Arbor, 1946.
\item
  A. WEIL, Foundations of algebraic geometry, Amer. Math. Soc. Coll. Public., vol.~XXIX, New York, 1946.
\end{enumerate}

